% Cross-domain example: LLM test-time optimisation -> online learning theory (dueling bandits).
% Query 10.48550_arxiv.2602.21585 (SIR-4 CS test, cross-field); gold q:4f02a72cb830786c
% ("Double Thompson Sampling for Dueling Bandits"). Numbers from colab_qualitative_cs.ipynb (8 Sep 2026):
%   rank of the gold: Qwen3 cosine 1,037 -> multi-view scorer 80 -> +OpenIE graph 10 -> +SciAffordGraph 32 -> +CCMP 36;
%   graph channel alone 1 (with and without CCMP); path weight 11.80 with CCMP, 12.06 without;
%   OpenIE graph channel 417, its top route is a five-hop co-mention chain of weight 0.01.
% Preamble needs: xcolor, tcolorbox[most], booktabs, float, tikz (positioning, arrows.meta, calc); \resulttablestyle.
% The colour/macro definitions are shared with fig_example_cs_rendering.tex (\providecommand, safe to \input both).
\providecolor{hlblue}{RGB}{214,228,247}
\providecolor{hlgreen}{RGB}{214,240,224}
\providecolor{hlorange}{RGB}{252,232,200}
\providecolor{hlpurple}{RGB}{228,220,246}
\providecolor{hlgrey}{RGB}{236,236,238}
\providecommand{\hl}[2]{\colorbox{#1}{\strut #2}}
\providecommand{\gw}[1]{{\color{black!55}(#1)}}
\providecommand{\ty}[1]{{\scriptsize\scshape\bfseries\color{black!75}#1}}
\providecommand{\qbox}[2]{%
  \begin{tcolorbox}[enhanced, colback=gray!4, colframe=black!70, boxrule=0.6pt, arc=2.5pt,
                    left=6pt, right=6pt, top=4pt, bottom=4pt, before skip=4pt, after skip=4pt]
  \textbf{#1}\quad #2
  \end{tcolorbox}}
\providecommand{\ladder}[5]{%
  {\footnotesize rank of this paper:\; Qwen3 cosine \textbf{#1} $\rightarrow$ multi-view scorer \textbf{#2}
  $\rightarrow$ {+}\,OpenIE graph \textbf{#3} $\rightarrow$ {+}\,SciAffordGraph \textbf{#4} $\rightarrow$ {+}\,CCMP \textbf{#5}}}

% =================================================================================================
\begin{figure}[H]
\centering
\small
\qbox{Research problem (Computer Science: language models)}{``How can we \hl{hlblue}{optimize candidates in a large
discrete output space} for test-time LLM improvement when \hl{hlorange}{informative scalar rewards are
unavailable or unreliable}? [\ldots] Some recent approaches avoid scalar evaluation entirely but only refine a
single trajectory against a single incumbent, lacking a global model of candidate quality and making
\hl{hlorange}{no principled allocation of a limited comparison budget}. No existing method
\hl{hlgreen}{aggregates many noisy pairwise preferences into a coherent global ranking} to drive search.''}

\qbox{Hypothetical answer $h_1$}{``A \hl{hlpurple}{multi-armed bandit framework that utilizes Thompson
sampling} to dynamically allocate resources for evaluating candidate solutions, optimizing the selection of
candidates based on \hl{hlgreen}{observed pairwise preferences}.''\quad {\footnotesize (best of six views;
cosine to the gold 0.67, raw query 0.39. View 2 names Copeland rank aggregation; views 3 to 5 are
cross-domain analogies from evolutionary biology, sociophysics and game theory.)}}

\qbox{Retrieved inspiration (Statistics: sequential decision theory)}{\emph{Double Thompson Sampling for Dueling Bandits}. A
bandit whose only feedback is a \hl{hlgreen}{noisy pairwise comparison} of two arms; D-TS keeps a posterior over
the preference matrix and draws both arms from it, with $O(K\log T)$ regret for Condorcet and Copeland
winners.\\[2pt]
\ladder{1{,}037}{80}{10}{32}{36}{\footnotesize; graph channel alone \textbf{1}.}}

\caption{Cross-domain example, SIR-4 CS, Computer Science to Statistics. A question about improving language-model outputs without a reward signal is
answered by dueling-bandit theory. Embedding similarity puts the paper at rank 1{,}037: the query talks
about language models and candidate refinement, the paper about arms and regret. The scorer's bandit view
lifts it to 80, and the SciAffordGraph reasoner ranks it first on its own channel through the query's
\textsc{limitation} frame ``inefficiency in candidate selection'' (Figure~\ref{fig:graph-cs-bandits}). The
fused rank stops at 36: the additive fusion under-weights a graph channel that is confident exactly where
the scorer is not, which is the fusion limitation discussed in the limitations section.}
\label{fig:qual-cs-bandits}
\end{figure}

% =================================================================================================
\begin{figure}[H]
\centering
\begin{tikzpicture}[
  font=\small, every node/.style={align=left},
  fr/.style={draw=black!70, line width=0.5pt, rounded corners=2.5pt, inner sep=5pt, text width=50mm},
  lim/.style={fr, fill=hlorange}, met/.style={fr, fill=hlpurple},
  ent/.style={fr, fill=hlgrey, text width=26mm, inner sep=3pt}, pap/.style={fr, fill=hlblue!60, text width=26mm, inner sep=3pt},
  qry/.style={fr, fill=white, dashed, text width=150mm, align=center},
  gold/.style={fr, fill=hlblue, line width=1.1pt, text width=60mm, align=center},
  e/.style={-{Latex[length=2.2mm]}, line width=0.7pt, color=black!80},
  ew/.style={-{Latex[length=2mm]}, line width=0.6pt, color=black!55},
  seed/.style={-{Latex[length=2mm]}, line width=0.6pt, color=black!45, dashed},
  rel/.style={font=\footnotesize\itshape, inner sep=2pt},
  relw/.style={font=\scriptsize\itshape, color=black!60, inner sep=1.5pt},
  note/.style={font=\scriptsize, color=black!60, text width=34mm},
  panel/.style={font=\small\bfseries, color=black!80},
]
\node[qry] (q) at (0,0) {\ty{query: Computer Science (language models)}\\ ``Optimise candidates in a large discrete space for test-time LLM improvement when scalar rewards are unavailable; aggregate noisy pairwise preferences into a global ranking''};
% ---- A: frame graph
\node[panel] at (-4.2,-1.8) {A. SciAffordGraph route (graph channel rank 1)};
\node[lim] (l1) at (-5.4,-3.1) {\ty{limitation}\\ inefficiency in candidate selection};
\node[met] (m1) at (-5.4,-5.7) {\ty{method}\\ double Thompson sampling (D-TS): draw both arms of each comparison from a posterior over the preference matrix};
\node[gold] (g) at (-4.6,-9.3) {\ty{gold paper: Statistics (sequential decision theory, dueling bandits)}\\ \emph{Double Thompson Sampling for Dueling Bandits}};
\draw[seed] (q.south) -- ++(0,-0.35) -| node[pos=0.3, above=3pt, font=\scriptsize, color=black!55] {seed frame of the query} (l1.north);
\draw[e] (l1) -- node[rel, right] {overcomes$^{-1}$} (m1);
\draw[e] (m1) -- node[rel, right] {contributes$^{-1}$} (g);
\node[note, anchor=north west] at ($(l1.north east)+(0.25,0)$) {the query's ``no principled allocation of a limited comparison budget'' seeds this frame};
\node[note, anchor=north west] at ($(m1.north east)+(0.25,0)$) {path weight 11.80 with CCMP \gw{12.06 without}; both hops gate 1.00};
% ---- B: OpenIE chain
\node[panel] at (5.4,-1.8) {B. OpenIE entity graph (graph channel rank 417)};
\node[ent] (o1) at (3.5,-3.1) {llm};
\node[ent] (o2) at (7.3,-3.1) {large language models};
\node[pap] (o3) at (7.3,-5.0) {\ty{paper}\\ Parallel-R1: parallel thinking via RL};
\node[ent] (o4) at (3.5,-5.0) {synthetic data};
\node[ent] (o5) at (3.5,-6.9) {d ts};
\draw[seed] (q.south) -- ++(0,-0.35) -| (o1.north);
\draw[ew] (o1) -- node[relw, above] {are$^{-1}$} (o2);
\draw[ew] (o2) -- node[relw, right] {is mentioned in} (o3);
\draw[ew] (o3) -- node[relw, below] {is mentioned in$^{-1}$} (o4);
\draw[ew] (o4) -- node[relw, right] {is based on$^{-1}$} (o5);
\node[gold, text width=60mm, fill=white, draw=black!50, dashed, line width=0.6pt] (g2) at (5.4,-9.3) {\ty{gold paper}\\ \emph{Double Thompson Sampling for Dueling Bandits}\\[1pt] {\footnotesize reached in five hops, path weight 0.01}};
\draw[ew] (o5.south) -- node[relw, right] {is mentioned in} (g2.north -| o5.south);
\end{tikzpicture}
\caption{Graph view of the dueling-bandits example. \textbf{A.} The SciAffordGraph reaches the paper in two
typed hops: the query's \textsc{limitation} on allocating comparisons is what D-TS overcomes, and D-TS is
what the paper contributes. \textbf{B.} The OpenIE entity graph reaches the same paper only through a
five-hop chain of co-mentions (``llm'' to ``large language models'' to an unrelated RL paper to ``synthetic
data'' to the token ``d ts''), with a path weight of 0.01 and a graph-channel rank of 417. Dashed arrows are
the seeding step; $r^{-1}$ is the inverse relation.}
\label{fig:graph-cs-bandits}
\end{figure}

% =================================================================================================
\begin{table}[H]
\centering
\resulttablestyle
\setlength{\tabcolsep}{5pt}
\renewcommand{\arraystretch}{1.25}
\begin{tabular}{@{}p{0.11\linewidth} p{0.87\linewidth}@{}}
\toprule
\multicolumn{2}{@{}l}{\textbf{Paths}\quad test-time optimisation without scalar rewards $\rightarrow$ \emph{Double Thompson Sampling for Dueling Bandits}} \\
\midrule
SciAfford\-Graph &
11.80 \gw{12.06}:\; (\textsc{limitation}: inefficiency in candidate selection,\; \textit{overcomes}$^{-1}$,\;
\textsc{method}: double Thompson sampling) $\rightarrow$
(\textsc{method}: double Thompson sampling,\; \textit{contributes}$^{-1}$,\; \textbf{gold paper}) \newline
{\footnotesize shortest seed-to-gold distance 2 hops; graph channel rank 1 with and without CCMP; fused rank 36 and 32.} \\
OpenIE &
0.01:\; (llm,\; \textit{are}$^{-1}$,\; large language models) $\rightarrow$ (large language models,\; \textit{is mentioned in},\;
\textsc{paper}: Parallel-R1) $\rightarrow$ (\textsc{paper}: Parallel-R1,\; \textit{is mentioned in}$^{-1}$,\; synthetic data) $\rightarrow$
(synthetic data,\; \textit{is based on}$^{-1}$,\; d ts) $\rightarrow$ (d ts,\; \textit{is mentioned in},\; \textbf{gold paper}) \newline
{\footnotesize shortest seed-to-gold distance 4 hops; graph channel rank 417; OpenIE model rank 10 (its scorer channel, rank 10).} \\
\bottomrule
\end{tabular}
\caption{Path interpretations for Figure~\ref{fig:qual-cs-bandits}. Top-weighted route from the query's seed
frames to the gold under each trained reasoner (beam search over the gradient of the graph score with
respect to each layer's edge weights, as in NBFNet and GFM-RAG); $r^{-1}$ is the inverse relation; grey
weight is the same weights with the CCMP gate off. The frame graph states the reason in two typed hops; the
entity graph's route is a chain of co-mentions with no content.}
\label{tab:paths-cs-bandits}
\end{table}
